% Compile in Overleaf with XeLaTeX.
\documentclass[a4paper,12pt]{article}

\usepackage{fontspec}
\usepackage{polyglossia}
\setmainlanguage{russian}
\setotherlanguage{english}
\setmainfont{DejaVu Serif}[
  BoldFont={DejaVu Serif Bold},
  ItalicFont={DejaVu Serif},
  BoldItalicFont={DejaVu Serif Bold}
]
\setsansfont{DejaVu Sans}
\setmonofont{DejaVu Sans Mono}

\usepackage{amsmath,amssymb}
\usepackage{geometry}
\usepackage{graphicx}
\usepackage{booktabs}
\usepackage{tabularx}
\usepackage{array}
\usepackage{float}
\usepackage{caption}
\usepackage{subcaption}
\usepackage{enumitem}
\usepackage{xcolor}
\usepackage{listings}
\usepackage{microtype}
\usepackage{fancyhdr}
\usepackage[hidelinks]{hyperref}

\geometry{left=25mm,right=20mm,top=20mm,bottom=22mm}
\setlength{\parindent}{1.25cm}
\setlength{\parskip}{0.15em}
\emergencystretch=6em
\renewcommand{\arraystretch}{1.18}
\captionsetup{font=small,labelfont=bf}
\setlist[itemize]{leftmargin=1.4cm,itemsep=0.2em}
\setlist[enumerate]{leftmargin=1.4cm,itemsep=0.2em}

\definecolor{codeblue}{RGB}{31,78,121}
\definecolor{codegreen}{RGB}{55,86,35}
\definecolor{codegray}{RGB}{245,245,245}
\lstset{
  language=Python,
  basicstyle=\ttfamily\footnotesize,
  keywordstyle=\color{codeblue}\bfseries,
  commentstyle=\color{codegreen},
  stringstyle=\color{purple},
  backgroundcolor=\color{codegray},
  frame=single,
  rulecolor=\color{gray!40},
  numbers=left,
  numberstyle=\tiny\color{gray},
  breaklines=true,
  showstringspaces=false,
  columns=fullflexible,
  keepspaces=true
}

\pagestyle{fancy}
\fancyhf{}
\fancyhead[L]{Алгоритмы компьютерной графики}
\fancyhead[R]{Лабораторная работа № 2}
\fancyfoot[C]{\thepage}
\setlength{\headheight}{15pt}

% ---------- Edit these fields before submission ----------
\newcommand{\studentname}{Чжун Ц. Су Л.}
\newcommand{\studentgroup}{указать группу}
\newcommand{\teachername}{указать преподавателя}
\newcommand{\reportyear}{2026}

\begin{document}

\begin{titlepage}
  \thispagestyle{empty}
  \begin{center}
    \textbf{МИНИСТЕРСТВО НАУКИ И ВЫСШЕГО ОБРАЗОВАНИЯ РФ}\\[0.6em]
    \textbf{УНИВЕРСИТЕТ ИТМО}\\[0.6em]
    Факультет программной инженерии и компьютерной техники

    \vfill

    {\Large\bfseries ОТЧЁТ ПО ЛАБОРАТОРНОЙ РАБОТЕ № 2}\\[1.2em]
    {\large по дисциплине «Алгоритмы компьютерной графики»}\\[2em]
    {\Large\bfseries Расчёт яркости на сфере\\
    от точечных источников света}

    \vfill

    \begin{tabular}{@{}ll@{}}
      Выполнил: & \studentname \\
      Группа: & \studentgroup \\
      Преподаватель: & \teachername
    \end{tabular}

    \vfill
    Санкт-Петербург\\
    \reportyear
  \end{center}
\end{titlepage}

\tableofcontents
\newpage

\section{Цель и постановка задачи}
Цель работы — освоить расчёт и визуализацию распределения энергетической яркости сферы, освещённой несколькими точечными источниками с ламбертовской диаграммой излучения. Отражение описывается диффузной и зеркальной составляющими модели Блинн–Фонга; зеркальная составляющая формирует выраженные блики.

Необходимо разработать приложение на Python, сформировать изображение заданного разрешения, нормировать яркость на её максимум в диапазон 0–255, сохранить изображение и вывести его на экран. К отчёту прилагаются исходный код, результирующее изображение, абсолютная яркость в трёх различных точках сферы, максимальное и минимальное значения.

Рекомендуемые диапазоны: W и H от 100 до 10000 мм; Wres и Hres от 200 до 800 пикселей с обязательным условием квадратных пикселей; xL, yL, xC, yC от \ensuremath{-}10000 до 10000 мм; zL и zC от 100 до 10000 мм; I0 от 0,01 до 10000 Вт/ср. Наблюдатель находится в O=(0,0,zO). Сфера должна целиком помещаться в поле зрения. Материал выбирается так, чтобы блик был хорошо заметен.

Используется правая мировая система координат. В демонстрационном варианте наблюдение направлено вдоль \ensuremath{-}Z, экран лежит в плоскости z=0, а оптические оси обоих источников направлены вдоль \ensuremath{-}Z, как на схеме задания. Программа также позволяет изменять направление наблюдения и расстояние от наблюдателя до экрана. Фоновое освещение не добавляется.

\subsection{Принятые единицы и модель отражения}
Сила излучения источников задана в Вт/ср. Поэтому рассчитывается энергетическая яркость L в Вт/(м\textsuperscript{2}·ср), а не фотометрическая яркость в кд/м\textsuperscript{2}. Перевод в фотометрические величины требует спектральных данных. В соответствии со схемой задания применяется эмпирическая функция отражения f=kd+ks·max(n·h,0)\textasciicircum{}m. Коэффициенты kd и ks трактуются как величины с размерностью ср\textsuperscript{-}\textsuperscript{1}. Модель не считается строго энергосохраняющей; дополнительный нормирующий множитель к зеркальной части не вводится.

\clearpage
\section{Теоретические сведения}
\subsection{Камера и перспективная проекция}
Пусть d — единичное направление наблюдения, er и eu — единичные направления вправо и вверх в плоскости экрана, f — расстояние от наблюдателя до экрана. Для каждого пикселя используется его центр; первая строка изображения соответствует верхнему краю экрана.

\begin{equation}
x_j=-W/2+(j+1/2)W/W_{\mathrm{res}}
\end{equation}
\begin{equation}
y_i=H/2-(i+1/2)H/H_{\mathrm{res}}
\end{equation}
\begin{equation}
\mathbf{q}_{ij}=\frac{f\mathbf{d}+x_j\mathbf{e}_r+y_i\mathbf{e}_u}{\|f\mathbf{d}+x_j\mathbf{e}_r+y_i\mathbf{e}_u\|}
\end{equation}
Направления экрана строятся с помощью векторных произведений. Если направление наблюдения почти параллельно исходному вектору вверх, выбирается другой вспомогательный вектор. Это предотвращает вырожденный базис камеры.

\subsection{Пересечение луча со сферой}
Луч задаётся уравнением P(t)=O+tq. Для центра C и радиуса R уравнение сферы имеет вид |P\ensuremath{-}C|\textsuperscript{2}=R\textsuperscript{2}. После подстановки получается квадратное уравнение; поскольку q единичный, коэффициент при t\textsuperscript{2} равен единице.

\begin{equation}
b=\mathbf{q}\cdot(\mathbf{O}-\mathbf{C}),\quad c=\|\mathbf{O}-\mathbf{C}\|^2-R^2
\end{equation}
\begin{equation}
t^2+2bt+c=0,\quad D=b^2-c,\quad t=-b-\sqrt{D}
\end{equation}
Если D<0, луч не пересекает сферу и пиксель остаётся фоновым. Наблюдатель проверяется на нахождение вне сферы; при D≥0 выбирается ближайшее положительное пересечение. Оно соответствует видимой поверхности и автоматически учитывает самозакрытие непрозрачной сферы.

\begin{equation}
\mathbf{P}=\mathbf{O}+t\mathbf{q},\qquad \mathbf{n}=\frac{\mathbf{P}-\mathbf{C}}{R}
\end{equation}
\clearpage
\subsection{Ламбертовские точечные источники}
Для источника Li вычисляются расстояние ri и единичное направление li от поверхности к источнику. В формуле освещения расстояние выражается в метрах, хотя геометрические входные данные заданы в миллиметрах.

\begin{equation}
r_i=\|\mathbf{L}_i-\mathbf{P}\|/1000,\qquad \mathbf{l}_i=\frac{\mathbf{L}_i-\mathbf{P}}{\|\mathbf{L}_i-\mathbf{P}\|}
\end{equation}
Единичная оптическая ось источника a=(0,0,\ensuremath{-}1), а луч от источника к поверхности направлен вдоль \ensuremath{-}li. Косинус угла излучения равен a·(\ensuremath{-}li)=li,z. Отрицательные значения обнуляются, поскольку источник не излучает в заднюю полусферу.

\begin{equation}
I_i=I_{0i}\max(l_{i,z},0)
\end{equation}
\begin{equation}
E_i(\mathbf{P})=\frac{I_{0i}}{r_i^2}\max(l_{i,z},0)\max(\mathbf{n}\cdot\mathbf{l}_i,0)
\end{equation}
В этой формуле первый косинус относится к диаграмме источника, второй — к проекции принимающей площадки. Они имеют разный смысл и должны учитываться независимо. При n·li≤0 точка находится на неосвещённой стороне относительно данного источника.

\subsection{Расчёт яркости по модели Блинн–Фонга}
\begin{equation}
\mathbf{v}=\frac{\mathbf{O}-\mathbf{P}}{\|\mathbf{O}-\mathbf{P}\|},\qquad \mathbf{h}_i=\frac{\mathbf{l}_i+\mathbf{v}}{\|\mathbf{l}_i+\mathbf{v}\|}
\end{equation}
\begin{equation}
f_i=k_d+k_s\max(\mathbf{n}\cdot\mathbf{h}_i,0)^m
\end{equation}
\begin{equation}
L(\mathbf{P},\mathbf{v})=\sum_i E_i(\mathbf{P})f_i
\end{equation}
Полная яркость складывается из диффузной части ΣEi·kd и зеркальной части ΣEi·ks·max(n·hi,0)\textasciicircum{}m. Увеличение m сужает блик. Вырожденный полувектор обрабатывается без деления на ноль. Если n·v≤0, направление на наблюдателя находится под поверхностью и выходная яркость принимается равной нулю.

\subsection{Нормировка изображения}
\begin{equation}
G_{ij}=\mathrm{round}(255L_{ij}/L_{\max})
\end{equation}
Lmax определяется только по видимым пикселям сферы. Фон имеет G=0 и не участвует в статистике. Если Lmax=0, формируется полностью чёрное изображение. Абсолютные значения сохраняются отдельно и не заменяются нормированными значениями серого.

\clearpage
\section{Исходные данные}
\begin{table}[H]
\centering
\caption{Параметры демонстрационного варианта}
\begin{tabular}{@{}p{0.51\textwidth}p{0.39\textwidth}@{}}
\toprule
Параметр & Значение \\
\midrule
Размер экрана W\ensuremath{\times}H & 1200\ensuremath{\times}1200 мм \\
Разрешение & 600\ensuremath{\times}600 пикселей \\
Наблюдатель O & (0, 0, 2500) мм \\
Направление наблюдения & (0, 0, \ensuremath{-}1) \\
Расстояние до экрана f & 2500 мм \\
Центр сферы C & (0, 0, 600) мм \\
Радиус R & 350 мм \\
Коэффициент kd & 0,20 ср\textsuperscript{-}\textsuperscript{1} \\
Коэффициент ks & 0,80 ср\textsuperscript{-}\textsuperscript{1} \\
Показатель блеска m & 80 \\
\bottomrule
\end{tabular}
\end{table}
\begin{table}[H]
\centering
\caption{Точечные источники света}
\begin{tabular}{@{}p{0.15\textwidth}p{0.40\textwidth}p{0.33\textwidth}@{}}
\toprule
Источник & Координаты, мм & I0, Вт/ср \\
\midrule
L1 & (\ensuremath{-}650, 450, 1700) & 100 \\
L2 & (550, \ensuremath{-}250, 1600) & 70 \\
\bottomrule
\end{tabular}
\end{table}
Размер пикселя по обеим осям равен 1200/600=2 мм. Условие квадратных пикселей выполнено. Расстояние от наблюдателя до центра сферы dC=1900 мм, поэтому наблюдатель находится вне сферы.

Для центрального положения сферы радиус её перспективного силуэта на экране равен fR/√(dC\textsuperscript{2}\ensuremath{-}R\textsuperscript{2})=468,544609 мм. Он меньше половины ширины и высоты экрана, равной 600 мм. Поэтому весь силуэт входит в изображение. Для смещённой сферы программа проверяет расстояния до всех четырёх боковых плоскостей поля зрения, а не только положение центра.

При kd=0,20, ks=0,80 и m=80 зеркальная составляющая образует два локальных пика. Положение каждого пика определяется полувектором света и наблюдения, поэтому оно не обязано совпадать с проекцией источника или центром сферы.

\clearpage
\section{Алгоритм приложения}
1. Ввод размеров и разрешения экрана, положения наблюдателя, направления наблюдения, расстояния до экрана, центра и радиуса сферы, параметров материала и списка источников.

2. Проверка конечности чисел, рекомендуемых диапазонов, квадратных пикселей и положительного радиуса. Источники должны лежать вне сферы; вся сфера должна находиться перед наблюдателем и внутри поля зрения.

3. Построение ортонормированного базиса камеры и координат центров пикселей. Верхняя строка соответствует положительному направлению вертикальной координаты экрана.

4. Формирование перспективных лучей и решение квадратного уравнения для каждой пары луч–сфера. Создание маски попаданий и координат ближайших пересечений.

5. Вычисление нормали, направления к наблюдателю, направлений к источникам, расстояний, косинусов излучения и падения. Затем вычисляются полувекторы и обе составляющие яркости.

6. Суммирование вкладов источников, вычисление минимумов и максимумов по видимым пикселям, нормировка в uint8 и сохранение PNG. Для контроля сохраняется отдельная маска сферы.

7. Прямой расчёт в трёх геометрически заданных точках. Дополнительно оцениваются экстремумы по угловой сетке всей поверхности: 361 полярный угол и 720 азимутов.

8. Сохранение цветной карты, графика сечения, CSV, JSON и NPZ с абсолютными яркостями, координатами поверхности и маской. Вывод изображения и результатов в графическом интерфейсе.

Основной расчёт векторизован с помощью NumPy. Для N=Wres·Hres пикселей и K источников сложность равна O(KN), память — O(N). Дополнительная сетка поверхности содержит 259920 отсчётов и оценивается отдельно; её максимум является дискретной оценкой, а не точным непрерывным экстремумом.

Программа lab2\_sphere.py запускается с интерфейсом без аргументов. Для воспроизводимого запуска используется команда python lab2\_sphere.py --no-gui --output results. Параметры можно передать через --config parameters.json. В ключевых местах реализации приведены пояснения на китайском языке.

\clearpage
\section{Результаты}
\subsection{Распределение яркости на сфере}
\begin{figure}[H]
\centering
\includegraphics[width=14cm]{results/brightness_map.png}
\caption{Распределение абсолютной энергетической яркости сферы}
\end{figure}
В изображении различимы два зеркальных блика и более широкая диффузная составляющая. У границы сферы часть поверхности не освещена обоими источниками, поэтому минимальная яркость видимой поверхности равна нулю. Белое поле на цветной карте показывает отсутствие пересечения луча со сферой; в итоговом восьмибитном PNG фон чёрный.

Результирующее нормированное изображение имеет размер 600\ensuremath{\times}600 пикселей, один восьмибитный канал, без гамма-коррекции и тонального преобразования. Значение 255 соответствует максимальной яркости текущего варианта. Карта физических значений и PNG используют одинаковое направление вертикальной оси.

\clearpage
\subsection{Нормированное изображение и сечение}
\begin{figure}[H]
\centering
\includegraphics[width=10cm]{results/sphere_normalized.png}
\caption{Восьмибитное изображение с нормировкой на максимум}
\end{figure}
\begin{figure}[H]
\centering
\includegraphics[width=14cm]{results/center_section.png}
\caption{Сечение через проекцию центра с разделением составляющих}
\end{figure}
Сечение вычислено отдельными лучами при yэкрана=0 и не выбирается из ближайшей строки изображения. В сечении показаны полная, диффузная и зеркальная яркости. Максимум сечения может быть ниже максимума изображения, поскольку оба блика смещены относительно горизонтальной оси.

\clearpage
\subsection{Три контрольные точки}
Контрольные точки задаются нормалями n1=(0,0,1), n2=(\ensuremath{-}0,6;0;0,8) и n3=(0,6;0;0,8). Координаты определяются как P=C+Rn. Все три точки находятся на сфере и видимы наблюдателю. Яркость вычисляется непосредственно по формулам, независимо от разрешения.

\begin{table}[H]
\centering
\caption{Абсолютные значения в контрольных точках}
\begin{tabular}{@{}p{0.15\textwidth}p{0.40\textwidth}p{0.33\textwidth}@{}}
\toprule
Точка & Координаты, мм & L, Вт/(м\textsuperscript{2}·ср) \\
\midrule
P1 & (0, 0, 950) & 17,674237 \\
P2 & (-210, 0, 880) & 14,117132 \\
P3 & (210, 0, 880) & 17,089102 \\
\bottomrule
\end{tabular}
\end{table}
\begin{table}[H]
\centering
\caption{Составляющие яркости, Вт/(м\textsuperscript{2}·ср)}
\begin{tabular}{@{}p{0.15\textwidth}p{0.40\textwidth}p{0.33\textwidth}@{}}
\toprule
Точка & Диффузная часть & Зеркальная часть \\
\midrule
P1 & 17,515758 & 0,158479 \\
P2 & 14,116729 & 0,000403 \\
P3 & 17,088220 & 0,000882 \\
\bottomrule
\end{tabular}
\end{table}
\subsection{Минимальная и максимальная яркость}
\begin{table}[H]
\centering
\caption{Статистика по двум различным способам дискретизации}
\begin{tabular}{@{}p{0.51\textwidth}p{0.39\textwidth}@{}}
\toprule
Характеристика & L, Вт/(м\textsuperscript{2}·ср) \\
\midrule
Минимум видимых пикселей & 0,000000 \\
Максимум видимых пикселей & 79,375576 \\
Среднее видимых пикселей & 14,944841 \\
Минимум по сетке всей сферы & 0,000000 \\
Максимум по сетке всей сферы & 79,325340 \\
\bottomrule
\end{tabular}
\end{table}
В изображении 172448 центров пикселей попадают на сферу. Максимум на изображении достигается около точки P=(94,574465; -43,216146; 934,197599) мм. Глобальная угловая сетка и сетка лучей не совпадают, поэтому их дискретные максимумы немного различаются. Нули фона не включены в минимум видимых пикселей. Нулевая яркость невидимой стороны учитывается только в дополнительной статистике всей сферы для направления на наблюдателя.

\clearpage
\section{Проверка корректности}
Выполнены восемь автоматических проверок в test\_lab2.py. Они проверяют аналитическое значение на оси, ближайшее пересечение луча со сферой, линейность по силе излучения, сложение независимых источников, нулевую яркость и невидимую сторону, принадлежность контрольных точек сфере, отклонение недопустимых параметров и корректность экспорта PNG. Все проверки завершились успешно.

Для независимой аналитической проверки оставлен один источник L=(0,0,1700) мм с I0=100 Вт/ср и точка P=(0,0,950) мм. В этой конфигурации cosθ=cosα=n·h=1, расстояние равно 0,75 м, kd+ks=1 ср\textsuperscript{-}\textsuperscript{1}. Поэтому ожидаемая яркость равна 100/0,75\textsuperscript{2}=177,777778 Вт/(м\textsuperscript{2}·ср), что совпадает с результатом программы.

При удвоении I0 обоих источников абсолютные яркости удваиваются, а PNG после нормировки на собственный максимум остаётся неизменным. Сумма двух отдельно рассчитанных источников совпадает с совместным расчётом. При kd=ks=0 получается чёрное изображение без NaN и деления на ноль. Центральный луч пересекает сферу в (0,0,950) мм; угловой луч экрана проходит мимо неё.

Повторное чтение PNG подтверждает один канал, заданный размер, значения и порядок строк. Для видимых пересечений выполнено |P\ensuremath{-}C|=R. Программа отклоняет неквадратные пиксели, нулевое направление наблюдения, обрезанную сферу и источник внутри сферы.

\section{Вывод}
Разработано приложение для расчёта яркости сферы от нескольких ламбертовских точечных источников. Реализованы перспективная камера, пересечения лучей с поверхностью, учёт диаграммы источников и угла падения, диффузная и зеркальная составляющие Блинн–Фонга, нормировка, экранная визуализация и сохранение результатов. Выбранные параметры дают два выраженных блика.

Для демонстрационного варианта яркость видимых пикселей находится в диапазоне от 0,000000 до 79,375576 Вт/(м\textsuperscript{2}·ср). Контрольные значения и проверки согласуются с принятой моделью. Полученные абсолютные значения относятся к выбранной эмпирической функции отражения; конечные сетки дают дискретные оценки экстремумов.

\section{Список источников}
1. Лабораторные работы по курсу «Алгоритмы компьютерной графики». ЛР 2 «Расчёт яркости на сфере от точечных источников света»: задание и схемы.

2. Материалы лекций по курсу «Алгоритмы компьютерной графики», предоставленные вместе с лабораторными работами.

3. Lagarde S. Adopting a physically based shading model. 2011. https://seblagarde.wordpress.com/2011/08/17/hello-world/.

\end{document}
